% Plan: development/outline.md, section 6, item G.
% Extended-value proofs of Lemma 4.1, Lemma 4.2 and Propositions 4.3-4.5 (as
% stated in sections/04-split.tex), and the proof of thm:lnts-opt with its
% lemma and proposition (moved here from Section 4.2 for length). The band theory (decision D1) is not part
% of the paper. Corrections of the pattern-theory critique applied: C3 (finite
% potentials, b_t > -infinity), C16 (F_t defined on all of R^{d_t}), C15 (the
% "only if" direction needs attainment). The proof of Lemma 4.1(c) covers
% arbitrary sets Z_c (via eq:split-subset). Round 2 (adjudication G2-01, moves
% M1 and M2): part (c) of the proof of prop:split-affine with the paragraph
% after it, and the "Best constants" part of the proof of prop:split-cellwise
% with the remark on the cap, moved to the supplement (B10-split-extras.tex,
% prop:split-forced and prop:split-cellwise-best).
\section{Proofs for the split certificates}\label{app:splitproofs}

We prove the validity results of \cref{sec:split-framework}, classical in substance, in the form the certificates use (extended-value costs, constrained stages, enclosures of unbounded states), and then \cref{thm:lnts-opt}.

\subsection{Setting and conventions}\label{app:splitproofs-setting}

We use the path model \eqref{eq:split-path}, the stage residuals $F^\varphi_t$
and split bounds $B(\varphi;K')$ of \eqref{eq:split-bound}, and enclosures
$K'_t\supseteq\proj_t\feas[\mathrm P]$, as defined in
\cref{sec:split-framework}.
Each separator map selects $k_t$ distinct coordinates.
The terms $\varphi_0(\pi^L_1z)$ and $\varphi_N(\pi^R_Nz)$ are read as $0$.

We use $\inf\emptyset=+\infty$, and a finite sum that contains the term
$+\infty$ equals $+\infty$, even if it also contains $-\infty$.
Two facts follow directly.
\begin{itemize}
\item[(F1)] For sets $A_1,\dots,A_m$ and functions
  $G_j\colon A_j\to\R\cup\{+\infty\}$,
  $\inf_{a\in A_1\times\dots\times A_m}\sum_j G_j(a_j)=\sum_j\inf_{A_j}G_j$.
  If some right-hand infimum is $+\infty$, both sides are $+\infty$.
  Otherwise every infimum lies in $[-\infty,+\infty)$, and near-minimizers of
  the separate terms give near-minimizers of the sum.
\item[(F2)] For $c\in\R$ and $G$ with values in $[-\infty,+\infty]$,
  $\inf(G+c)=\inf G+c$.
\end{itemize}

\subsection{Split bounds}\label{app:splitproofs-bound}

\begin{proof}[Proof of \cref{lem:split-bound}]
We first prove a slightly stronger form of (a).
Let $S\subseteq\feas[\mathrm P]$ and let $K'_t\supseteq\proj_tS$ for every
$t$.
Then
\begin{equation}\label{eq:split-subset}
  B(\varphi;K')\le\inf\Bigl\{\sum_{t=1}^NF_t(z_t):z\in S\Bigr\}.
\end{equation}
If $S=\emptyset$, the right side is $+\infty$.
Let $z\in S$ and write $s_t=\pi^R_tz_t=\pi^L_{t+1}z_{t+1}$ for $1\le t<N$.
Each value $\varphi_t(s_t)$ enters $F^\varphi_t(z_t)$ with a plus sign and
$F^\varphi_{t+1}(z_{t+1})$ with a minus sign, at the same point $s_t$.
All terms are finite, so
\begin{equation}\label{eq:split-telescope}
  \sum_{t=1}^N F^\varphi_t(z_t)
  =\sum_{t=1}^N F_t(z_t)+\sum_{t=1}^{N-1}\bigl[\varphi_t(s_t)-\varphi_t(s_t)\bigr]
  =\sum_{t=1}^N F_t(z_t).
\end{equation}
Since $z_t\in\proj_tS\subseteq K'_t$, every stage infimum satisfies
$\inf_{K'_t}F^\varphi_t\le F^\varphi_t(z_t)<+\infty$.
So $B(\varphi;K')$ is a sum of terms in $[-\infty,+\infty)$, and
$B(\varphi;K')\le\sum_tF_t(z_t)$; the infimum over $z\in S$ gives
\eqref{eq:split-subset}.

(a) is \eqref{eq:split-subset} with $S=\feas[\mathrm P]$.

(b) The proof of \eqref{eq:split-subset} shows that a feasible point makes
every stage infimum, and hence $B(\varphi;K')$, smaller than $+\infty$.
If every $F_t$ vanishes on $K_t\cap\operatorname{dom}F_t$, every feasible
point has objective $0$, so $\vstar\in\{0,+\infty\}$, and by (a)
$B(\varphi;K')>0$ forces $\vstar=+\infty$.

(c) If $\feas[\mathrm P]=\emptyset$, then $\vstar=+\infty$.
Otherwise every $z\in\feas[\mathrm P]$ lies in some $Z_c$, so
$\sum_tF_t(z_t)\ge\beta_c\ge\min_c\beta_c$; take the infimum over $z$.
For the last sentence, apply \eqref{eq:split-subset} with
$S=\feas[\mathrm P]\cap Z_c$: for every split $\varphi^c$ and all sets
$K^c_t\supseteq\proj_t(\feas[\mathrm P]\cap Z_c)$, the number
$\beta_c=B(\varphi^c;K^c)$ is valid.
If $Z_c$ is a box, that is, a product of intervals for some model variables,
such sets are obtained by intersecting enclosures $K'_t$ with the intervals for
the variables that bag $t$ contains.

(d) For $z\in\feas[\mathrm P]\subseteq Y$ the added terms vanish, so the
infimum over $Y$ is at most $\sum_tF_t(z_t)$; take the infimum over $z$.
The infeasibility statement follows as in (b).
With $\mu=(\lambda_1,\dots,\lambda_{N-1})$ and the copy rows, the added terms
are $\sum_{t<N}\lambda_t^\top(\pi^R_tz_t-\pi^L_{t+1}z_{t+1})$, and with
$Y=K'_1\times\dots\times K'_N$ the infimum separates by (F1) into the stage
infima of $F_t(z)+\lambda_t^\top\pi^R_tz-\lambda_{t-1}^\top\pi^L_tz$; this is
$B(\varphi;K')$ for the affine split with slopes $\lambda$
(\cref{prop:split-affine}(a)).
\end{proof}

The proof of (a) uses only that each $\varphi_t$ is real-valued, which the
telescoping in \eqref{eq:split-telescope} needs.

\subsection{Value-function minorants}\label{app:splitproofs-minorant}

We use the setting of \cref{lem:split-minorant}: states $y_i\in Y_i$, controls
$u_i\in U_i$ (nonempty) and maps $g_i\colon Y_{i-1}\times U_i\to Y_i$ for
$1\le i\le n$, a terminal cost $\Phi\colon Y_n\to\R$, and
$\vstar=\inf\{\Phi(y_n):y_0\in Y_0,\ u_i\in U_i,\ y_i=g_i(y_{i-1},u_i)\}$.
Stage costs can be included by adding the accumulated cost to the state.
The value functions are $V_n=\Phi$ and
$V_{i-1}(y)=\inf_{u\in U_i}V_i(g_i(y,u))$ for $y\in Y_{i-1}$ and
$i=n,\dots,1$.

\begin{lemma}[dynamic programming]\label{lem:split-dp}
For every $i$ and $y\in Y_i$, $V_i(y)$ is the infimum of $\Phi(y_n)$ over all
controls $u_{i+1},\dots,u_n$, where $y_i=y$ and $y_j=g_j(y_{j-1},u_j)$.
In particular $\vstar=\inf_{y\in Y_0}V_0(y)$.
\end{lemma}

\begin{proof}
Backward induction on $i$; the case $i=n$ is the definition.
For the step, the infimum over $(u_i,\dots,u_n)$ equals the infimum over $u_i$
of the infimum over the remaining controls, which is
$V_i(g_i(y,u_i))$ by the induction hypothesis.
\end{proof}

\begin{proof}[Proof of \cref{lem:split-minorant}]
(a) $\inf_{Y_0}W_0\le\inf_{Y_0}V_0=\vstar$ by \cref{lem:split-dp}.

(b) First, $w_j\le V_{i-1}(r_j)$: for every $u$, $M_i(u)r_j\in\R^2_+$, so
$W_i(M_i(u)r_j)\le V_i(M_i(u)r_j)$, and the infimum over $u$ is
$V_{i-1}(r_j)$ by the definition of $V_{i-1}$.
Second, $V=V_{i-1}$ is superadditive.
Positive homogeneity gives $V(0)=\alpha V(0)$ for all $\alpha>0$, so $V(0)=0$
and $V(\alpha y)=\alpha V(y)$ for all $\alpha\ge0$; with concavity,
$V(y+y')=2V\bigl(\tfrac12y+\tfrac12y'\bigr)\ge V(y)+V(y')$.
Hence, for $\sigma,\sigma'\ge0$,
\[
  V(\sigma r_j+\sigma'r_{j+1})\ge\sigma V(r_j)+\sigma'V(r_{j+1})
  \ge\sigma w_j+\sigma'w_{j+1}.
\]
Third, the interpolant is well defined.
The angles of the $r_j$ increase strictly from $0$ to $\pi/2$, so the sectors
$\operatorname{cone}\{r_j,r_{j+1}\}$ cover $\R^2_+$, consecutive sectors meet
along a common ray and other pairs only at $0$.
A point inside a sector has a unique representation
$\sigma r_j+\sigma'r_{j+1}$, a point on a common ray gets the same value from
both sectors, and $W_{i-1}(0)=0$.
\end{proof}

\Cref{lem:catmix-dp} applies (b) stage by stage to the \inst{catmix} grids.

\subsection{Affine splits, windows and cellwise slopes}\label{app:splitproofs-affine}

\begin{proof}[Proof of \cref{prop:split-affine}]
(a) By (F2), the constant $c_t-c_{t-1}$ can be taken out of the stage infimum
of bag $t$, and these constants sum to $c_N-c_0=0$ (with $c_0=c_N=0$); if a
stage infimum is $+\infty$, the bound is $+\infty$ for every $c$.
The remaining stage residuals are
$F_t(z)+\lambda_t^\top\pi^R_tz-\lambda_{t-1}^\top\pi^L_tz$.
The Lagrangian of the copy formulation,
$\sum_tF_t(z_t)+\sum_{t<N}\lambda_t^\top(\pi^R_tz_t-\pi^L_{t+1}z_{t+1})$, is the
sum of these functions over the bags, so by (F1) its infimum over
$K_1\times\dots\times K_N$ equals the sum of their infima.
The same argument applies to $B(\varphi;K')$ with the sets $K'_t$.

(b) The argument does not use affinity.
If $\bar z$ minimizes every stage residual, then by
\eqref{eq:split-telescope} and \cref{lem:split-bound}(a)
\[
  B(\varphi)=\sum_tF^\varphi_t(\bar z_t)=\sum_tF_t(\bar z_t)\ge\vstar\ge B(\varphi).
\]
Conversely, if $\vstar=B(\varphi)$ is attained at the feasible point $z^*$,
then $\vstar$ is finite and
$\sum_tF^\varphi_t(z^*_t)=\vstar=\sum_t\inf_{K_t}F^\varphi_t$, so every stage
infimum is finite.
Each difference $F^\varphi_t(z^*_t)-\inf_{K_t}F^\varphi_t$ is nonnegative and
their sum is $0$, so each is $0$.
\end{proof}

\begin{proof}[Proof of \cref{prop:split-window}]
Let $W'_{ab}$ be the set over which $\beta_{[a,b]}(\varphi;K')$ is taken.
On $W'_{ab}$ the split terms at the separators $a,\dots,b-1$ telescope as in
\eqref{eq:split-telescope} (where an $F_t$ is $+\infty$, both sides are
$+\infty$), so the window objective equals $\sum_{t=a}^bF^\varphi_t(z_t)$.
Its infimum over $W'_{ab}$ is at least its infimum over
$K'_a\times\dots\times K'_b$, which is $\sum_{t=a}^b\inf_{K'_t}F^\varphi_t$ by
(F1).
This gives the first inequality; both sides are $+\infty$ if an outside term
is.
For the second, merge bags $a,\dots,b$ of (P) into one bag with variable
$(z_a,\dots,z_b)$, the set of $(z_a,\dots,z_b)\in K_a\times\dots\times K_b$
whose consecutive copies agree, cost $\sum_{t=a}^bF_t$, left separator map
$\pi^L_a$ and right separator map $\pi^R_b$.
The merged path model has the same feasible points, grouped differently, and
the same value $\vstar$.
The set $W'_{ab}$ contains the merged bag's components of all feasible points,
so it is an enclosure for the merged bag, and the split
$(\varphi_1,\dots,\varphi_{a-1},\varphi_b,\dots,\varphi_{N-1})$ has merged
stage infimum $\beta_{[a,b]}(\varphi;K')$.
\Cref{lem:split-bound}(a) for the merged model gives the second inequality.
For several disjoint windows, merge each of them.
\end{proof}

\begin{proof}[Proof of \cref{prop:split-cellwise}]
If $\feas[\mathrm P]=\emptyset$ there is nothing to prove.
Let $z\in\feas[\mathrm P]$, $s_t=\pi^R_tz_t=\pi^L_{t+1}z_{t+1}$, and choose a
cell $D_t\in\mathcal P_t$ that contains $s_t$ (any one, if there are several;
both bags at separator $t$ then use the chosen cell).
The term $\lambda_{t,D_t}^\top s_t$ appears with a plus sign in bag $t$ and
with a minus sign in bag $t+1$, so
\[
  \sum_tF_t(z_t)=\sum_{t=1}^N\bigl[F_t(z_t)+\lambda_{t,D_t}^\top s_t
  -\lambda_{t-1,D_{t-1}}^\top s_{t-1}\bigr],
\]
with the terms for $t=0$ and $t=N$ read as $0$.
Bag $t$ satisfies $z_t\in K'_t$, $\pi^L_tz_t\in D_{t-1}$ and
$\pi^R_tz_t\in D_t$, so its bracket is at least the infimum that bounds
$b_t(D_{t-1},D_t)$ from above in \cref{prop:split-cellwise}, and hence at
least $b_t(D_{t-1},D_t)$.
The brackets are finite, so these $b_t$ are finite, the path $(D_t)$ has
length at most $\sum_tF_t(z_t)$, and $\mathrm{SP}\le\sum_tF_t(z_t)$.
\end{proof}

\subsection{Proof of \texorpdfstring{\cref{thm:lnts-opt}}{the lnts theorem}}\label{app:splitproofs-lnts}

We use the notation of \cref{sec:split-affine}: $w_j$, $c_j$, $r_j$, $\alpha(h)$, $\beta(h)$, and $C$, $D$ and $g$ of \eqref{eq:lnts-g}.

\begin{lemma}[exact elimination]\label{lem:lnts-elim}
For every $\theta\in\R^{N+1}$ and $h\in\R$, the dynamics rows and the zero initial states determine all states, and $v^x_N=100h\sum_jw_j\cos\theta_j$, $v^y_N=100h\sum_jw_j\sin\theta_j$ and $p^y_N=100h^2\sum_jc_j\sin\theta_j$.
Hence a point of \inst{lnts}$N$ is exactly feasible if and only if its states follow from the dynamics, $|\theta_j|\le\bar\theta$ for all $j$, $h>0$, and
\begin{equation}\label{eq:lnts-terminal}
  \sum_jw_j\cos\theta_j=\alpha(h),\qquad \sum_jw_j\sin\theta_j=0,\qquad \sum_jc_j\sin\theta_j=\beta(h).
\end{equation}
\end{lemma}

\begin{proof}
Each row has coefficient~1 on its next state.
Summing the velocity rows, and then the vertical position rows, gives the three expressions (\cref{app:lnts-proof} computes the weights).
The value $p^x_N$ is free, $h=0$ gives $v^x_N=0\ne45$, and for $h>0$ the terminal conditions become \eqref{eq:lnts-terminal}.
\end{proof}

\begin{proposition}[support bound]\label{prop:lnts-support}
For $\mu\in\R$ and $\nu\ge0$ put $S(\mu,\nu)=\sum_j\bigl(w_j^2+(\mu w_j+\nu c_j)^2\bigr)^{1/2}$.
Every exactly feasible point of \inst{lnts}$N$ satisfies $\alpha(h)+\nu\beta(h)\le S(\mu,\nu)$.
\end{proposition}

\begin{proof}
By \eqref{eq:lnts-terminal}, $\alpha(h)+\nu\beta(h)=\sum_j[w_j\cos\theta_j+(\mu w_j+\nu c_j)\sin\theta_j]$; apply the Cauchy--Schwarz inequality to each term.
\end{proof}

\begin{proof}[Proof of \cref{thm:lnts-opt}]
\emph{The root.}
We have $g(0)=-20N^2/81<0$.
As $\nu\to\infty$, $C(\nu)\to1$, because only $j=N/2$ has $r_j=0$, and $D(\nu)\to\sum_jw_j|r_j|\ge(N-1)/2$, so $g$ has a positive limit.
On $\nu\ge0$, $D'>0$ and $C'\le0<C$, so $g'=D'-\frac{40}{81}CC'>0$, and $g$ has exactly one positive root.

\emph{Feasibility.}
Put $t_j=\nu^*r_j$, so that $\cos\theta^*_j=(1+t_j^2)^{-1/2}$ and $\sin\theta^*_j=t_j(1+t_j^2)^{-1/2}$; the angle bounds hold because $|\theta^*_j|<\pi/2<\bar\theta$.
Since $w_{N-j}=w_j$ and $r_{N-j}=-r_j$ (\cref{lem:lnts-weights}), the sines cancel in pairs, and $\sum_jw_j\sin\theta^*_j=0$.
Next, $\sum_jw_j\cos\theta^*_j=C(\nu^*)=\alpha(h^*)$, and, with $c_j=w_j(r_j+N/2)$, $\sum_jc_j\sin\theta^*_j=\sum_jw_jr_j\sin\theta^*_j=D(\nu^*)=\frac{20}{81}C(\nu^*)^2=\beta(h^*)$.
By \cref{lem:lnts-elim} the point is exactly feasible, with objective $Nh^*$.

\emph{Optimality.}
Take $\mu^*=-\nu^*N/2$ in \cref{prop:lnts-support}.
Then $\mu^*w_j+\nu^*c_j=w_jt_j$, and by the three identities of the feasibility step $S(\mu^*,\nu^*)=\sum_jw_j\sqrt{1+t_j^2}=\sum_j[w_j\cos\theta^*_j+w_jt_j\sin\theta^*_j]=\alpha(h^*)+\nu^*\beta(h^*)$.
So every exactly feasible point has $\alpha(h)+\nu^*\beta(h)\le\alpha(h^*)+\nu^*\beta(h^*)$.
The left side is strictly decreasing in $h>0$, so $h\ge h^*$ and $Nh\ge Nh^*$.

\emph{Enclosure.}
Since $C$ is nonincreasing, rationals $0<a<b$ with $g(a)<0<g(b)$ give $9N/(20\,C(a))\le\vstar\le9N/(20\,C(b))$.
The certificate of \cref{thm:lnts-opt} (\cref{app:lnts-proof}) finds such $a,b$ with $b-a<10^{-90}$, and outward rational bounds on $C(a)$ and $C(b)$ give an enclosure of width below \sci{1.11}{-91}, which we round outward to 16 decimals.
\end{proof}
